\clearpage
\VerdierPhysicalPageAnchor{0233}%
\begingroup\small\setlength{\parskip}{0pt}\setlength{\emergencystretch}{1.5em}
\setlength{\abovedisplayskip}{4pt plus 2pt minus 2pt}%
\setlength{\belowdisplayskip}{4pt plus 2pt minus 2pt}%

\VerdierSemanticAnchor{ch4d-4-1-8}
\noindent\textbf{4.1.8.}\quad
当 $X$ 是单复形时，由按第一次数滤过的复形 $X$ 所定义的、取值于
$K(\mathcal A)$ 的谱对象 (4.1.7)，正是与复形 $X$ 相联系的、取值于
$K(\mathcal A)$ 的第一谱对象。

\VerdierSemanticAnchor{ch4d-4-1-8-degree-support}
类似地，设 $X$ 是 $n$-重复形，$m$ 是整数，$1\leq m\leq n$。假设 $X$
的分量 $X^\sigma$ 在下列条件成立时为零：
\[
\sigma-\sigma_m e_m\neq0.
\]
这里，$\sigma_m$ 表示第 $m$ 个坐标。于是可验证，由按第 $m$ 个次数滤过的
复形 $X$ 所定义的谱对象等于单复形 $\int_m X$ 的第一典范谱对象。

\VerdierSemanticAnchor{ch4d-4-1-8-derived-image}
最后，当范畴 $\mathcal A$ 是阿贝尔范畴时，典范函子
\[
K(\mathcal A)\longrightarrow D(\mathcal A)
\]
把由按第 $m$ 个次数滤过的 $n$-重复形 $X$ 所定义的、取值于
$K(\mathcal A)$ 的谱对象变为取值于 $D(\mathcal A)$ 的谱对象。如此得到的
谱对象称为\emph{由按第 $m$ 个次数滤过的 $X$ 所定义的、取值于
$D(\mathcal A)$ 的谱对象}。不会导致混淆时，省略“取值于 $D(\mathcal A)$”
这句话。

\VerdierSemanticAnchor{ch4d-4-1-9}
\medskip\noindent\textbf{4.1.9.}\quad
设 $X^{\bullet\bullet}$ 是 $\mathcal A$ 中的双复形。

\VerdierSemanticAnchor{ch4d-4-1-9-row-complexes}
对每个整数 $q$，以 $X^{q,\bullet}$（相应地，$X^{\bullet,q}$）表示单复形
\[
\begin{aligned}
X^{q,\bullet}
&=\cdots\xrightarrow{d_2^{q,i-1}}X^{q,i}
  \xrightarrow{d_2^{q,i}}X^{q,i+1}
  \xrightarrow{d_2^{q,i+1}}\cdots,\\[3pt]
(\text{resp.}\quad X^{\bullet,q}
&=\cdots\longrightarrow X^{i,q}
  \xrightarrow{(-1)^q d_1^{i,q}}X^{i+1,q}
  \xrightarrow{(-1)^q d_1^{i+1,q}}X^{i+2,q}
  \longrightarrow\cdots\quad),
\end{aligned}
\]
其次数 $i$ 的分量为 $X^{q,i}$（相应地，$X^{i,q}$）。

\VerdierSemanticAnchor{ch4d-4-1-9-direction-morphisms}
$X^{\bullet\bullet}$ 在第一方向（相应地，第二方向）上的微分诱导态射
（在 $K(\mathcal A)$ 中）：
\[
\begin{aligned}
d_1^{q,\bullet}&\colon X^{q,\bullet}\longrightarrow X^{q+1,\bullet}\\
(\text{resp.}\quad
d_2^{\bullet,q}&\colon X^{\bullet,q}\longrightarrow X^{\bullet,q+1}\quad),
\end{aligned}
\]

\VerdierSemanticAnchor{ch4d-4-1-9-components}
其分量为
\[
\begin{aligned}
(d_1^{q,\bullet})^i&=d_1^{q,i}\\
(\text{resp.}\quad(d_2^{\bullet,q})^i&=(-1)^i d_2^{i,q}\quad).
\end{aligned}
\]
\endgroup
